% introduction_overview.tex
Every factor analysis of a rating scale begins with items written in ordinary
language, and we take that wording seriously as data. A scale's items constitute a
\emph{semantic proposal}---a deliberate attempt, in language, to sample the
content that constitutes a construct---and modern language-model embeddings make
that proposal measurable before any respondent answers, mapping each item to a
point in a vector space whose geometry encodes its meaning. We introduce
\texttt{semanticfa} \citep{yanitski-westbury-2026}, an R package for
\emph{semantic factor analysis} (SFA): exploratory factor analysis of the
similarity structure of language-model embeddings of scale items, requiring no
human response data at any stage.

The motivating distinction is between two relational structures over the same
items. Classical factor analysis recovers a scale's \emph{behavioral} structure,
the latent dimensions of person-by-item response covariance; SFA recovers its
\emph{semantic} structure, the geometry of item meaning. These are structurally
parallel analyses on different data sources, and their agreement,
\emph{semantic--behavioral coherence}, is an estimable relation in a construct's
nomological network: convergence is intranetwork validity evidence, while
divergence is itself informative, locating where the linguistic and behavioral
renderings of a construct come apart. Throughout, SFA \emph{complements}, never
replaces, response-based factor analysis. We develop this in three movements---the
theory of analyzing a scale's language, a review of language models in
psychometrics, and the positioning of \texttt{semanticfa} and its demonstration.
